Passive forensic detectors answer whether an image was edited, but they supply no path back to what the image showed before the edit. This paper presents a deterministic, training-free, key-seeded self-embedding fragile watermark that localizes tampering and restores content jointly. Each $8{\times}8$ block carries, in its two least-significant bit-planes, a 32-bit truncated HMAC-SHA256 tag over its own six most-significant bit-planes and a 96-bit DCT recovery descriptor of a spatially distant partner block selected by a key-seeded single-cycle permutation. A flagged block is rebuilt from the descriptor its surviving partner carries, or explicitly marked unrecoverable. The evaluation covers a 32-image corpus in 3,232 runs over four tamper classes and seven tamper ratios, plus a worked example on an author-supplied $512{\times}512$ photograph. Watermarked images reach 45.48~dB with SSIM 0.9885. Block recall is 1.0000, with block precision structurally forced to 1.0000 by the ground-truth rule, and pixel precision is 0.9599. The recoverability rate $\rho$ stays above its structural bound $1-\alpha$ at every ratio, with mean 0.6482, and in-region recovery reaches 28.49~dB for Variant~A. The scheme is fully fragile, and its one-to-one backup allocation fails all-or-nothing per block, so reference-sharing schemes degrade more gracefully.
